\begin{claimboundary}
OpenAIの2月stackは、Codex app、GPT-5.3-Codex、System Card、Sparkを一つのartifactに潰さない。workspace、base model、安全性文書、realtime-oriented research previewは役割が異なる。Sparkの1,000 tokens/s超という値はCerebras hardwareとresearch-preview条件に依存するvendor claimであり、一般的なCodex throughputとして扱わない。System Cardのcyber classificationも、記載されたthreshold/precautionの境界を越えて強く言い換えない。\autocite{src-15823ae5556e,src-4a841ea7ec7c,src-9aba7995b6a5,src-dd64fe3084a4}
\end{claimboundary}

\begin{claimboundary}
Qwen3.5/Qwen3-Coder-Nextのchronologyは地域差を保持する。Model StudioのChinaとinternationalではavailabilityとdated entryが異なり、living lifecycle pageには後日のglobal rolloutが追記され得る。したがって2月号では2月16日・20日の記録を2月の事実として固定し、後から見えるcurrent availabilityを遡及して使わない。\autocite{src-eac8dfbeb78c}
\end{claimboundary}

\begin{claimboundary}
Voxtral Realtimeについて確定できるのは、2月4日のrelease、streaming speech-to-text、open weights、Apache 2.0、4B、13 languagesといった範囲である。sub-200ms latencyやaccuracy比較はserving configurationに依存するvendor claimであり、open-weightであることも周辺hosted service全体のopen-source化を意味しない。またEvidenceはMistral news indexを起点としているため、item-level表現をFebruary launch articleの記載から広げない。\autocite{src-d8b7c1c36830}
\end{claimboundary}

\begin{claimboundary}
Deep ThinkはGemini 3 family全体へ一般化できる能力ではなく、science/research/engineering向けのspecialized reasoning modeとして扱う。2月発表時点のAPI accessはlimited/early accessであり、科学benchmarkや比較結果もvendor-reportedである。後日のavailabilityや結果を2月へ遡及させず、2月時点で確認できたspecializationの方向だけを残す。\autocite{src-dd5d9b8662f3}
\end{claimboundary}

\begin{claimboundary}
vLLM v0.16.0は2月25日releaseだが、branch cutは2月8日である。このためlate-Februaryのmodel updateがすべて同releaseへ統合済みだったとは言えない。またseriesにはv0.15.1とv0.16.0が含まれるが、systems signalの中心はv0.16.0である。release noteのthroughput/TPOT改善値もproject measurementとして扱い、独立benchmarkへ昇格させない。\autocite{src-5041feb6f87f,src-67307025f5d9}
\end{claimboundary}

この総括で意図的に残したいのは、『確認できた変化』と『まだ強く言えないこと』の差である。複数agent workspace、repository-level tool interaction、streaming ASR、specialized reasoning、Realtime/tool serving APIは2月の一次資料で確認できる。一方、benchmark superiority、latency・throughput、科学性能などの数値比較はvendor/project claimのままである。2月が後続月へ残したのは、勝者の順位ではなく、AIを継続実行可能なsystemにするための論点がmodel・training・interaction・runtime・safetyへ分解され、同時進行し始めたという技術地図である。\autocite{src-15823ae5556e,src-4a841ea7ec7c,src-5041feb6f87f,src-67307025f5d9,src-9aba7995b6a5,src-d8b7c1c36830,src-dd5d9b8662f3,src-dd64fe3084a4,src-eac8dfbeb78c}


したがって3月以降との接続は、後続月のニュースを2月へ逆輸入する形では行わない。後続月を読むときに問うべきなのは、2月に立ち上がった五つの軸が継続・変形・収束したかである。agentはより長く安全に仕事を持てるようになったか。runtimeは新しいmodalityとtool contractを吸収したか。open modelはtrainingからdeploymentまで一貫したstackを作れたか。specialized reasoningは汎用modelと異なる価値を定着させたか。2月号は、その後の変化を比較可能にする基準点として閉じる。\autocite{src-5041feb6f87f,src-67307025f5d9,src-9aba7995b6a5,src-d8b7c1c36830,src-dd5d9b8662f3,src-eac8dfbeb78c}
